\section{Discussion}
The results align with the design intent: each optimization produces its largest effect on the class of scenes for which it was designed, and the optimizations as a group target the lower-capability portion of the hardware spectrum. For educational deployments, where institutions cannot dictate which devices students use, this targeting pattern is precisely what is needed. The desktop-tier result -- that the optimizations produce no measurable change because the hardware was already sufficient -- is not a negative finding but a confirmation that the engine imposes no overhead on capable hardware while delivering substantial benefit on constrained hardware.

\subsection{Why frame stability matters more than peak FPS in education}
The most important finding of the evaluation is that the optimizations primarily affect frame-time variability, not average performance. A scene rendering at an average of 110 FPS but with occasional 37 ms spikes feels less responsive to a student than a scene rendering at 120 FPS with consistent 15 ms maxima. For learning content, where students may interact with the same scene over extended sessions, perceived smoothness affects engagement more than peak throughput. The 97.5\% reduction in maximum frame time variability is therefore the most practically significant result of this work.

\subsection{Limitations of the current engine}
Three limitations are worth noting. First, the engine assumes Three.js as the rendering backend; the adapter layer is not yet generalized, although the architecture admits such generalization. Second, advanced features common in commercial engines -- entity-component-system architecture, GPU-driven culling, level-of-detail systems, OffscreenCanvas-based rendering -- are not yet implemented. Third, the engine does not yet provide a graphical scenario authoring tool; scenarios are authored as TypeScript code, which limits accessibility for educators without programming expertise.